\documentclass[a4paper,12pt]{article}
\usepackage[utf8]{inputenc}
\usepackage[spanish]{babel}
\usepackage{amsmath}
\usepackage{amsfonts}
\usepackage{cite}
\usepackage{graphicx}
\usepackage{wrapfig}
\usepackage{array}
\usepackage{listings}
\usepackage{tcolorbox}
\usepackage{enumerate}
\usepackage{enumitem}


\begin{document}

% --- Portada personalizada ---
\begin{titlepage}
    \centering
    \vspace*{1cm}

    {\Huge \textbf{Práctica 2}} \\[0.5cm]
    % {\Huge \textbf{Práctica 1}} \\[0.2cm]
    % {\large \textit{Sección práctica}} \\[0.5cm]
    {\Large {David García Curbelo}}

    \vspace*{1cm}
    \rule{0.05\linewidth}{0.3mm} \\[0.4cm]
    \vspace*{1cm}

    % Representación de Textos en Espacios Vectoriales y Probabilísticos \\[0.2cm] # para textos largos
    \begin{large}
        Minería de Datos \\[0.2cm]
        Universidad Nacional de Educación a Distancia \\[0.5cm]
        Curso 2024-2025
    \end{large}

    \vfill

    \begin{minipage}{0.8\textwidth}
        \centering
        \vspace{0.5cm}
        \rule{0.9\linewidth}{0.3mm} \\[0.4cm]
        Master Universitario en Tecnologías del Lenguaje
    \end{minipage}
\end{titlepage}

\newpage
\tableofcontents
\newpage


% --- Contenido ---


\section{Generación de datos}


El conjunto de datos sintético generado para esta práctica está compuesto por cinco variables predictoras y una variable objetivo (clase), con un total de 100 instancias. De las 5 variables predictoras, las variables \(X_1, X_2, X_3\) y \(X_5\) se definen como variables booleanas generadas aleatoriamente con una distribución uniforme. La variable \(X_4\) se define como la operación XOR bit a bit entre las variables \(X_2\) y \(X_3\):

\[
    X_4 = X_2 \oplus X_3
\]

\noindent donde \( \oplus \) representa la operación XOR (OR exclusivo). Esta función devuelve un valor de 1 si el número de variables entre \( X_2 \) y \( X_3 \) con valor 1 es impar, y 0 si es par.

La variable objetivo, denotada como \( Y \), se define como una función XOR extendida a tres dimensiones:

\[
    Y = X_1 \oplus X_2 \oplus X_3,
\]

\noindent Esta definición garantiza que la relación entre las variables predictoras y la clase sea no lineal y conjunta, de modo que ninguna variable individual sea suficiente para predecir \( Y \) de manera aislada.

Desde el punto de vista de la relevancia de las variables, \(X_1\) se considera relevante en sentido fuerte, ya que su exclusión modificaría la probabilidad de predecir correctamente la variable objetivo. Las variables \(X_2\) y \(X_3\) son relevantes en sentido débil, pues aunque su contribución a la predicción de \(Y\) es importante, su información puede ser parcialmente recuperada a través de \(X_4\), que depende exclusivamente de ellas. La variable \(X_4\), a pesar de derivarse de \(X_2\) y \(X_3\), sigue siendo relevante en sentido débil dado que introduce redundancia estructurada en el conjunto de datos y puede influir en los métodos de selección de características. Finalmente, \(X_5\) es una variable completamente irrelevante, ya que su valor es generado de manera completamente aleatoria e independiente del resto de las variables.

El diseño de este conjunto de datos permite evaluar con precisión las técnicas de selección de características, al contar con una relevancia conocida de antemano para cada variable. Además, la inclusión de una variable irrelevante y otra redundante permite analizar la robustez de los algoritmos ante el ruido y situaciones de selección de características complejas. Este marco sintético proporcionado en \cite{kohavi1997wrappers} nos genera un entorno controlado para comparar distintas estrategias de selección de características en algoritmos de aprendizaje supervisado, garantizando una evaluación rigurosa de su desempeño y de sus posibles limitaciones.



\newpage
\section{Técnicas}

Las técnicas de selección de características permiten identificar un subconjunto óptimo de variables que contribuyan de manera significativa a la predicción de la variable objetivo, eliminando aquellas redundantes o irrelevantes. Para el conjunto de datos sintetizado, se aplicarán cinco técnicas de selección de características, incluyendo tres métodos de filtrado, el análisis de componentes principales (PCA) y un método de envoltura basado en un clasificador supervisado.

\subsection{Métodos de filtrado}

\subsubsection{Información mutua}

Los métodos de filtrado evalúan la relevancia de las características de manera independiente del modelo de clasificación, aplicando criterios estadísticos o teóricos para determinar su importancia. Uno de los enfoques más comunes dentro de este grupo es la selección basada en la información mutua, donde se mide la dependencia estadística entre cada variable predictora y la variable objetivo. Formalmente, la información mutua entre dos variables aleatorias \( X \) e \( Y \) se define como:

\[
    I(X, Y) = \sum_{x \in X} \sum_{y \in Y} p(x,y) \log \frac{p(x,y)}{p(x)p(y)}
\]

\noindent donde \( p(x,y) \) representa la distribución conjunta de \( X \) e \( Y \), y donde \( p(x) \) y \( p(y) \) son las distribuciones marginales. Un valor alto de información mutua indica que el conocimiento de \( X \) reduce significativamente la incertidumbre sobre \( Y \), sugiriendo su utilidad en la clasificación.

\subsubsection{Test de independencia \(\chi^2\)}

Otro método dentro de los filtros es el test de independencia de \(\chi^2\), empleado para evaluar si una variable predictora y la variable objetivo son independientes. Se basa en la comparación entre las frecuencias observadas y las frecuencias esperadas bajo la hipótesis de independencia, calculando el estadístico:

\[
    \chi^2 = \sum_{i} \sum_{j} \frac{(O_{ij} - E_{ij})^2}{E_{ij}}
\]

\noindent donde \( O_{ij} \) y \( E_{ij} \) representan, respectivamente, la frecuencia observada y la esperada para la combinación de valores \( i, j \). Un valor alto de este estadístico sugiere que la variable en cuestión está fuertemente asociada con la clase.

\subsubsection{Relief}

El tercer método de filtrado a emplear es Relief, un algoritmo basado en la evaluación de la capacidad discriminativa de cada variable en relación con las instancias más cercanas en el espacio de características. Relief asigna un peso a cada variable en función de su capacidad para diferenciar correctamente ejemplos de diferentes clases, ajustándolo iterativamente a partir de comparaciones con los vecinos más próximos de la misma clase (``near-hit") y de clases opuestas (``near-miss"). Este método resulta especialmente útil en problemas donde las interacciones entre características juegan un papel clave, ya que captura dependencias locales entre variables.

\subsection{PCA}

El análisis de componentes principales constituye una técnica ampliamente utilizada para la reducción de dimensionalidad y la selección de características, aunque no por ello óptima en la mayoría de situaciones. Su aplicación difiere de los métodos anteriores, ya que transforma las variables originales en nuevas combinaciones lineales denominadas componentes principales. Estos componentes se obtienen mediante la descomposición en valores singulares de la matriz de covarianza de los datos, generando una base ortogonal donde las nuevas variables son ordenadas según la varianza explicada. Dado un conjunto de características \( X \), se busca una transformación lineal \( W \) tal que:

\[
    Z = WX
\]

\noindent donde las filas de \( W \) corresponden a los vectores propios de la matriz de covarianza de \( X \), ordenados en función de sus valores propios. Seleccionando los primeros \( k \) componentes con mayor varianza, se obtiene una representación compacta de los datos que conserva la mayor cantidad de información posible.

\subsection{Método de envoltura}

Finalmente, el método de selección basado en envoltura se fundamenta en la evaluación iterativa de subconjuntos de características mediante la optimización del rendimiento de un clasificador. En este caso, se empleará un enfoque de búsqueda heurística combinada con validación cruzada para identificar el subconjunto de variables que maximiza la precisión del modelo. Este método presenta la ventaja de adaptarse específicamente a la naturaleza del clasificador utilizado, evitando posibles incompatibilidades entre la selección de características y la fase de modelado. Sin embargo, su principal inconveniente radica en el coste computacional, ya que requiere entrenar múltiples modelos para evaluar diferentes subconjuntos de variables.


\newpage
\section{Propiedades}


\newpage
\section{Resultados}


\newpage
\section{Conclusiones}


\newpage
\addcontentsline{toc}{section}{Referencias}
\bibliographystyle{plain}
\bibliography{library}

\end{document}
